% #############################################################################
% Abstract Text
% !TEX root = ../main.tex
% #############################################################################
% reset acronyms
\acresetall
% use \noindent in firts paragraph
\noindent Session-typed programming languages based on linear logic, such as CLASS, enforce memory safety and deadlock-freedom by construction. The Linear Session Abstract Machine (SAM) was introduced to provide a formal model for sequentially executing the processes of session-typed languages, allowing for orders of magnitude better runtime performance than existing interpreters which spawn threads for each concurrent process. However, the SAM remains an interpreter, which introduces a performance overhead compared to a compiled language, even though the language is very strongly statically typed. This project introduces the first intermediate representation (IR) for session-based languages, which can be used to compile session-typed programs to efficient low-level imperative code. A compilation scheme for CLASS, which translates the high-level session-typed language into this IR, while maintaining all of the safety and liveness guarantees of the source language, was developed. An implementation of this scheme has been developed which covers all of the language features and which implements multiple optimization techniques. It is now possible to compile and run programs written in CLASS, achieving performance orders of magnitude better than the SAM interpreter.